\chapter{Conclusions and Future Work}
\label{chap:conclusion}

This thesis addressed the gap between linguistically plausible coordination and justified executable coordination in heterogeneous autonomous systems. With respect to the first objective, the state-of-the-art analysis showed that classical multi-agent systems, symbolic AI, LLM agents, and interoperability protocols contribute necessary but individually insufficient mechanisms. Their unresolved intersection is the explicit authorization boundary between language-mediated proposals and coordinated action.

With respect to the second objective, the thesis defined a protocol-independent framework in which capabilities, artifacts, dependencies, authority, commitments, validators, and completion evidence are explicit. Its central result is a total decision at the planning boundary: the coordinator produces either an authorized plan supported by predicate evidence or a structured infeasibility report. With respect to the third objective, that framework was mapped to a Python \texttt{CoordinationAgent} design over \texttt{CoordinationSdk} middleware, with protocol adaptation, registry admission, capability indexing, local and linguistic agent runtimes, feasibility analysis, task monitoring, and artifact aggregation handled at the implementation layer. The mapping illustrates implementability without making any particular communication protocol, linguistic runtime, or local runtime part of the theoretical claim.

With respect to the fourth objective, the thesis executed a reproducible prototype evaluation based on feasible and infeasible deterministic scenarios, Docker A2A system checks, local LLM reference batches, and separate outcome, process, communication-boundary, coordination-ownership, and trace metrics. The evidence shows that the implementation can authorize feasible plans, refuse missing or unauthorized requirements before dispatch, separate runtime failure from planning infeasibility, preserve SDK-mediated communication, and record model-specific linguistic reference behavior. The thesis does not claim production readiness, guaranteed runtime success, or broad empirical superiority over all possible baselines. This boundary is itself consistent with the proposed framework: conclusions should be authorized by available evidence rather than by plausible expectation.

The principal contribution is therefore a precise integration pattern for controlled openness: language supports interpretation and proposal; symbols support authorization and traceability; independently deployed agents retain execution autonomy; and unsupported requests are refused explicitly. The framework answers RQ1--RQ3 at the level of design, formal specification, and prototype behavior. RQ4 is answered for the implemented evidence set by the reported outcome, trace, boundary, latency, and local-model metrics, while larger baseline comparisons remain future work.

Future work should extend the implemented \texttt{CoordinationAgent}, \texttt{CoordinationSdk}, deterministic scenario runner, JSONL ledger, Docker A2A harness, and local LLM reference runner into a larger paired baseline and ablation study. Subsequent research should evaluate dynamic registry changes, adversarial capability advertisements, cryptographic identity and attestation, replicated coordinators with lease-based hot standby, domain-specific semantic validators, human authorization for high-impact actions, and third-party A2A implementations. Evaluation across protocols, local non-linguistic agents, Lingo-authored linguistic agents, larger local and hosted language models, and real operational settings is necessary to determine which findings generalize beyond the first agent-over-SDK instantiation.
